\section*{Three exact finite models}\addcontentsline{toc}{section}{Three exact finite models}
\subsection*{A seed remembers multiplicity support}
For $M_2(\mathbb C)\otimes I_3$, a coefficient matrix with rank one generates an orbit of dimension two. A matrix of rank two generates an orbit of dimension four. Neither single seed spans the full space of dimension six, because a $2\times3$ matrix cannot have rank three. Two seeds can nevertheless cover all three multiplicity coordinates together. The example illustrates the finite seed rank formula with $d=2$ and $m=3$.
\begin{figure}[H]\centering
\includegraphics[width=.94\textwidth]{monograph_figures/exact_seed_support.pdf}
\caption{The figure displays exact coefficient matrices for $d=2$, $m=3$. The combined row support of two seeds has rank three. The bar heights record algebraic orbit dimensions and should not be read as probabilities of unitary state preparation.}\label{monograph:fig:seed_support}
\end{figure}
\clearpage
\subsection*{An eliminated coordinate still contributes}
Take $H=\left(\begin{smallmatrix}2&1\\1&3\end{smallmatrix}\right)$ and retain its first coordinate. The eliminated block is invertible and the Schur complement of $H$ is $2-1/3=5/3$. Its inverse is $3/5$, exactly the retained entry of $H^{-1}$. Inverting the compressed entry alone gives $1/2$. The difference records the contribution of a path through the eliminated coordinate. Direct compression omits that contribution, while block inversion retains it.
\begin{figure}[H]\centering
\includegraphics[width=.94\textwidth]{monograph_figures/exact_feshbach_compression.pdf}
\caption{This rational model uses two coordinates to illustrate block inversion. Both the eliminated block and the full operator are invertible, and the plotted values are exact fractions. The example shows why direct compression and resolvent elimination give different results.}\label{monograph:fig:feshbach_compression}
\end{figure}
\clearpage
\subsection*{Normalization changes the map}
Let $P=|0\rangle\langle0|$ on $\mathbb C^2$ and $Z=\operatorname{diag}(1,-1)$. The normalized Kraus family $\sqrt{3/4}\,I$, $\sqrt{1/4}\,Z$ satisfies $\sum_j K_j^*K_j=I$ and maps $P$ to $P$. The unnormalized family $\sqrt3\,I$, $Z$ satisfies $\sum_j E_j^*E_j=4I$ and maps $P$ to $4P$. The two families have the same relative weights, but their total normalization differs. Both outputs are proportional to the same projector, but their traces differ.
\begin{figure}[H]\centering
\includegraphics[width=.94\textwidth]{monograph_figures/exact_channel_scale.pdf}
\caption{The figure shows exact weights for one supported code state. The normalized map preserves trace one, while the unnormalized positive map has output trace four. The illustration records this finite scaling distinction and does not assert a general recovery construction.}\label{monograph:fig:channel_scale}
\end{figure}
\clearpage
